% Quantitative receivers for the selected trace chart.
\subsection{Weighted bound control for the trace receiver}
\label{sec:ex-lattice-trace-weighted-estimates}

This section constructs the normalized response estimates and their scalar
combination, which feeds the rank-one receiver lemma
\ref{lem:ex-lattice-trace-weighted-price}. The chosen and reference fields use the same finite-volume response columns,
arch and separation ultrafilter. We keep the field decay separate from
bound convergence: the former uses the priced mission and its detector.

\begin{lemma}[Weighted bound from a rank-one approximate kernel]
\label{lem:ex-lattice-trace-weighted-price}
Let $E$ and $H$ be finite-dimensional real normed spaces. Let
$f:E\to H$ and $U:H\to E$ be linear, with $fU=\mathrm{id}_H$ and
$\ker f=\mathbb R x$ for some $x\ne0$. Suppose linear maps
$v_i:E\to H$ and $u_i:H\to E$ satisfy
\[
 v_i\longrightarrow f,\qquad u_i\longrightarrow U,
 \qquad v_i u_i=\mathrm{id}_H\quad\text{eventually}.
\]
Let $w_i\to x$, let $e_i\in\mathbb R$, and let $\rho:E\to\mathbb R$
be a fixed linear functional. If
\[
 e_i v_i(w_i)\longrightarrow0\quad\text{in }H,
 \qquad e_i\rho(w_i)\longrightarrow0,
\]
then there are linear functionals $\mu_i:H\to\mathbb R$ with uniformly
bounded operator norms and
\[
 \bigl\|e_i(\mu_i v_i-\rho)\bigr\|_{E^*}\longrightarrow0.
\]
\end{lemma}
\begin{proof}
Choose a linear functional $\lambda:E\to\mathbb R$ with $\lambda(x)=1$,
by extending $x$ to a basis. Define
\[
 A(y)=\lambda\bigl(y-Uf(y)\bigr).
\]
Since $f(y-Uf(y))=0$, the kernel assumption gives
$y-Uf(y)=a x$ for some scalar $a$. Applying $\lambda$ shows $a=A(y)$,
and therefore
\begin{equation}
 y=Uf(y)+A(y)x. \label{eq:trace-weighted-estimates-display-star}
\end{equation}
Set $V_i=u_i f+A\otimes w_i$, where
$(A\otimes w_i)(y)=A(y)w_i$. The convergence of $u_i$ and $w_i$ implies
$V_i-\mathrm{id}_E=(u_i-U)f+A\otimes(w_i-x)$, by (*). The right side
tends to zero in operator norm. Once $\|V_i-\mathrm{id}_E\|<1$, the
Neumann series gives an inverse and
\[
 \|V_i^{-1}-\mathrm{id}_E\|
 \leq \frac{\|V_i-\mathrm{id}_E\|}{1-\|V_i-\mathrm{id}_E\|},
\]
so $V_i^{-1}\to\mathrm{id}_E$.

Put $\mu_i=\rho u_i$. These functionals have bounded norms because
$u_i\to U$. For $y\in E$, the exact identity $v_i u_i=\mathrm{id}_H$
gives
\begin{align*}
 e_i(\mu_i v_i-\rho)(V_i y)
 &=e_i\rho\bigl(u_i v_i(u_i f(y)+A(y)w_i)-u_i f(y)-A(y)w_i\bigr)\\
 &=A(y)\,e_i\bigl(\rho(u_i v_i(w_i))-\rho(w_i)\bigr).
\end{align*}
The first term inside the final difference vanishes after weighting,
because
\[
 e_i\rho(u_i v_i(w_i))=\rho(u_i(e_i v_i(w_i)))
\]
and $\rho u_i$ is uniformly bounded. The second term vanishes by hypothesis. Thus the
left side tends to zero for every fixed $y$. In finite dimension,
convergence on a basis implies operator-norm convergence, so
$e_i(\mu_i v_i-\rho)V_i\to0$. Multiplication by
$V_i^{-1}\to\mathrm{id}_E$ gives
$e_i(\mu_i v_i-\rho)\to0$ in operator norm.
\end{proof}

Let $E$ be a finite-dimensional real normed space, $H$ a finite-dimensional
real normed space, $s_i:E\to H$ linear maps, and $\rho:E\to\mathbb R$ a
fixed linear functional. Let $x\in E$ satisfy $\rho(x)=0$. For each $i$,
let $b_i\ge1$ with $b_i\to\infty$, let $q_i,r_i\in E$, and let
$e_i\ge0$. Define
\[
 m_i=b_i^{-1}q_i-r_i.
\]
Assume
\begin{align}
 b_i^{-1}q_i&\longrightarrow0,
 &\sup_i|\rho(q_i)|&<\infty,
 &r_i&\longrightarrow -x,
 \label{eq:ex-lattice-trace-response-limit}\\
 e_i/b_i&\longrightarrow0,
 &e_i\rho(r_i+x)&\longrightarrow0,
 &b_i s_i(m_i)&\longrightarrow0\quad\text{in }H.
 \label{eq:ex-lattice-trace-response-weighted}
\end{align}

\begin{lemma}[The medium has both weighted null limits]
\label{lem:ex-lattice-trace-medium-receipts}
Under \eqref{eq:ex-lattice-trace-response-limit}--
\eqref{eq:ex-lattice-trace-response-weighted},
\[
 m_i\longrightarrow x,\qquad
 e_i\rho(m_i)\longrightarrow0,\qquad
 e_i s_i(m_i)\longrightarrow0.
\]
\end{lemma}
\begin{proof}
Since $b_i^{-1}q_i\to0$ and $r_i\to-x$,
$m_i=b_i^{-1}q_i-r_i\to x$. Since $\rho(x)=0$,
\[
 e_i\rho(m_i)=(e_i/b_i)\rho(q_i)-e_i\rho(r_i+x).
\]
The first term tends to zero by boundedness of $\rho(q_i)$ and
$e_i/b_i\to0$, and the second by hypothesis. Also $e_i/b_i\le1$
eventually, so
\[
 \|e_i s_i(m_i)\|=(e_i/b_i)\|b_i s_i(m_i)\|
 \le\|b_i s_i(m_i)\|\longrightarrow0.
\]
\end{proof}

All sequences in the application below are already reindexed by a
selected level $k_i\to\infty$: $b_i$ is the
arch scale at level $k_i$, $q_i$ and $r_i$ are the chosen and reference
responses there, and $s_i$ is the chart at index $i$. Precisely, $k_i$ is
the summit index of the frame level of the $i$th selected chart. It is the
level called the postselected level and denoted $n_i$ in
Subsection~\ref{sec:trace-prosper-sensor}; it is not the initial level
denoted $k_i$ there. This reindexing matters:
the field-decay estimate is assumed, and used, only at the reindexed
levels, not at every level. Write $p_i=\widehat\beta_{k_i}$ and
$A_a(i)=\exp(a(\log\log p_i)^2)$. The weight used is
$e_i=A_{1/5}(i)$ at this selected level. The next subsection proves
$e_i/b_i\to0$ and
$e_i\le p_i$. The field input is supplied with weight $p_i$; since
$b_i\le2p_i$, this gives $b_i s_i(m_i)\to0$ as required above.

In the selected construction, $E$ is the real signed-equivariant space of
grade-five coefficient arrays, $x$ is its nonzero equation-of-motion
direction, and $\rho$ is the free scalar bound. The vectors $q_i$ and
$r_i$ are respectively the transported chosen response and the transported
reference response. Their combination $m_i=b_i^{-1}q_i-r_i$ is the
medium used to build a response functional. The scale $b_i$ is the positive
arch scale along the selected chart. Boundedness of $\rho(q_i)$ follows from the unweighted Ward-bound limit;
the weighted estimate on $\rho(r_i+x)$ is the scale-resolved
reference-bound comparison; and the last estimate is the field decay of
the medium.

Combining Lemma~\ref{lem:ex-lattice-trace-medium-receipts} with
Lemma~\ref{lem:ex-lattice-trace-weighted-price}, under the chart hypotheses
$s_i\to f$, $u_i\to U$, $fU=\mathrm{id}$, $\ker f=\mathbb Rx$ with
$x\ne0$, and $s_i u_i=\mathrm{id}$ eventually, gives bounded linear
functionals $\mu_i:H\to\mathbb R$ and
\[
 \sup_i\|\mu_i\|<\infty,
 \qquad
 \bigl\|e_i(\mu_i s_i-\rho)\bigr\|_{E^*}\longrightarrow0.
\]
This is the response-estimate used by the two-direction trace
correction. The displayed deduction is complete as a deduction from its
hypotheses. In the selected construction it is carried out in
Proposition~\ref{prop:trace-prosper-sensor}: that proposition takes the three limits of
Lemma~\ref{lem:ex-lattice-trace-medium-receipts} as its marked-vector
hypotheses, derives $s_i\to f$ from them together with its own chart
inputs, and then applies
Lemma~\ref{lem:ex-lattice-trace-weighted-price}. The convergence
$s_i\to f$ is therefore not a consequence of the present paragraph, and
the chart inputs of that proposition are distinct from the response
estimates proved below. The field estimate uses the separately priced
mission, with the exact same-medium identification at the end of this section.

\subsubsection{The rate and medium calculation}
\label{subsec:ex-lattice-trace-source-rate}

For the already selected levels $k_i\to\infty$, write
$p_i=\widehat\beta_{k_i}$ for the running positive rate and $b_i$ for the
arch scale at $k_i$. The arch is the fixed-level ultrafilter limit of the running couplings
$\beta_{N,\nu,k}$. On the ultrafilter-large set $N\ge k$, the
selected Wilson history gives
$\widehat\beta_k/2\le\beta_{N,\nu,k}\le2\widehat\beta_k$.
The interval is closed, so its ultrafilter limit $b(k)$ belongs to it.
In particular $b(k)>0$. The clock tends to infinity; substitution
$k=k_i\to\infty$ proves the scale comparison
\begin{equation}
 \frac12p_i\le b_i\le 2p_i,\qquad p_i\longrightarrow\infty.
 \label{eq:ex-lattice-trace-arch-comparison}
\end{equation}
The age weights are $A_a(i)=\exp(a(\log\log p_i)^2)$. They satisfy
\begin{equation}
 A_a(i)A_c(i)=A_{a+c}(i),\qquad
 A_a(i)\le p_i^\delta\quad\text{eventually for every fixed }a\in\mathbb R
 \text{ and every }\delta>0.
 \label{eq:ex-lattice-trace-age-rules}
\end{equation}
For the second relation, take logarithms. Since
$(\log\log p)^2=o(\log p)$, eventually
$|a|(\log\log p_i)^2\le\delta\log p_i$. Exponentiation gives the bound.
Also $A_a(i)\to0$ directly for $a<0$, since
$(\log\log p_i)^2\to\infty$.

In particular, for $a\in\{1/6,1/5\}$,
\begin{equation}
 0\le \frac{A_a(i)}{b_i}
 \le 2p_i^{a-1}\longrightarrow0.
 \label{eq:ex-lattice-trace-age-over-arch}
\end{equation}
The same argument gives $A_a(i)\le p_i$ eventually, so an estimate with
weight $p_i$ implies the corresponding estimate with weight $A_{1/5}(i)$.

Here is the scalar deduction from the response and reference receipts. Let
$q_i$ be the transported chosen response and $w_i$ the transported reference response. The local response estimate and the
separated flux limit have the form
\begin{equation}
 \|q_i\|\le C A_{1/6}(i),\qquad
 \rho(q_i)\longrightarrow-\kappa,
 \label{eq:ex-lattice-trace-response-inputs-source}
\end{equation}
while the reference chart and the reference-bound comparison give
\begin{equation}
 w_i\longrightarrow -x,\qquad \rho(x)=0,\qquad
 A_{1/5}(i)\rho(w_i)\longrightarrow0.
 \label{eq:ex-lattice-trace-reference-inputs-source}
\end{equation}
Define $m_i=b_i^{-1}q_i-w_i$. Then
\[
 \|b_i^{-1}q_i\|\le 2C p_i^{-5/6}\longrightarrow0,
\]
so $m_i\to x$. Since $\rho(q_i)$ converges it is bounded, and
\[
 A_{1/5}(i)\rho(m_i)
 =\frac{A_{1/5}(i)}{b_i}\rho(q_i)
   -A_{1/5}(i)\rho(w_i)\longrightarrow0
\]
by \eqref{eq:ex-lattice-trace-age-over-arch} and
\eqref{eq:ex-lattice-trace-reference-inputs-source}. Finally, if the local
field estimate is $p_i\,s_i(m_i)\to0$, then
$A_{1/5}(i)\le p_i$ gives $A_{1/5}(i)s_i(m_i)\to0$. These are exactly
the three limits consumed by the rank-one approximate-kernel lemma; the
remaining bound-receiver construction is its finite-dimensional argument,
not another analytic estimate.

The response and reference estimates are proved below from the local
charged coefficient and the paid response column of the same finite-volume
history. In particular the arch comparison is a closed-interval limit,
and the transport error is the normalized bias in Subsection~\ref{sec:ex-lattice-trace-coffer-supplier};
neither requires an additional rate assumption on the arch.

\subsubsection{How the two weighted receipts are assembled}
\label{subsec:ex-lattice-trace-weighted-assembly}

The weighted bound receipt for the medium uses only the response-bound
limit, the reference-bound cancellation, and the age-rate comparison just
proved. It is therefore useful to distinguish it from the stronger
estimate that supplies the weighted field image. If $\rho(q_i)\to-\kappa$,
then $|\rho(q_i)|\le M$ eventually for some $M$; thus
\[
 \left|\frac{A_{1/5}(i)}{b_i}\rho(q_i)\right|
 \le 2M p_i^{-4/5}\longrightarrow0.
\]
No weighted convergence rate for the response flux itself is needed here.
The separate reference-chart estimate is the cancellation
\[
                         A_{1/5}(i)\rho(w_i)\to0.
\] This identifies the precise roles of the two
limits and avoids conflating a field-size estimate with a bound
estimate.

For the field image, the construction proves the stronger statement
$p_i s_i(m_i)\to0$ in the finite-dimensional chart norm. The transfer to
the weaker age weight is the monotonicity estimate
$A_{1/5}(i)\le p_i$; there is no loss from summing over a basis because the
chart is a fixed finite-dimensional operator. The bound
$\|q_i\|=O(A_{1/6}(i))$ alone proves only $m_i\to x$ and does not imply this
field estimate. The latter is supplied by the local pricing argument and
its finite-volume passage.

\subsubsection{The response limits and the remaining estimates}
\label{subsec:ex-lattice-trace-weighted-route}

The response-size and response-flux limits above follow from a concrete
separation argument. Let $\mathcal T$ be the linear map from the transported
finite-volume coefficient field to the grade-five response, and let
$\lambda$ be the flux functional on the transported fields. On the limiting
field, their compatibility is $\rho(\mathcal T v_i)=\lambda(v_i)$. At fixed
level $i$, write $v_i(\nu)$ for the finite-volume transported field of the
chosen response, indexed by the finite-volume index $\nu$ (a different
object from the rate $p_i$). The local charged-configuration estimate
gives, on a set in the separation ultrafilter,
\begin{align}
 \|\mathcal T v_i(\nu)\|&\le C_TA_{1/6}(i),\\
 |\lambda(v_i(\nu))+\kappa|&\le
 C_1A_{-1/4}(i)+C_2A_{-1}(i)
       +\frac{p_i^{-30}+p_i^{-100}}{12}.
 \label{eq:ex-lattice-trace-finite-response-price}
\end{align}
Here $C_T,C_1,C_2$ are fixed constants from the finite-dimensional
coefficient, response, and flux maps. The first line is the field-cost part
of the estimate; the second is the flux-bound part.
The transported field converges along that ultrafilter to the infinite
volume field $v_i$. Since $\mathcal T$ and $\lambda$ are continuous linear
maps, the inequalities pass to their limits. Therefore
$q_i=\mathcal T v_i$ satisfies the response norm bound and
$\rho(q_i)=\lambda(v_i)\to-\kappa$. Every term on the right tends to zero by
\eqref{eq:ex-lattice-trace-age-rules}, proving
\eqref{eq:ex-lattice-trace-response-inputs-source}.

Here is the finite-volume derivation, including its normalization. Work
in the fixed coefficient frame, with its supremum norm, and let $P$ delete
the unit coordinate. Thus $\|Pz\|\le\|z\|$. The map $\mathcal T$ is
signed averaging followed by the fixed identification with real covariant
arrays. Put $\lambda=\rho\mathcal T$. Since the unit has zero vector
flux, $\lambda P=\lambda$. These are identities of these maps, not a
norm assertion about the nonlinear local coefficient construction.

At the fixed level $k_i$ write $\vartheta_i$ for the budget threshold,
$X_i(\nu)$ for the chosen coefficient array, and $B_i(\nu)$ for the
public labels produced from its atomic certificate. Its charge is
$q=\vartheta_i^2 A_{1/6}(i)$, and its attainment cost is at most $p_i$.
Apply the finite-volume response record specialization of
Subsection~\ref{sec:ex-lattice-trace-coffer-supplier}, with $u=1$, separately
to the four entries. It gives the actual columns $\mathcal C_{k_i,\nu}$
of those labels and
\begin{equation}
 \|X_i(\nu)-\mathcal C_{k_i,\nu}B_i(\nu)\|_\infty
 \le \vartheta_i^2 A_{-1/4}(i).
 \label{eq:trace-weighted-paid-bias}
\end{equation}
Indeed its column error is at most
$4(1+c_{k_i})\vartheta_i^2A_{1/6}(i)/A_{1/2}(i)$, and
$4(1+c_{k_i})\le A_{1/12}(i)$ eventually. The exponent is exactly
$1/12+1/6-1/2=-1/4$. The tail is chosen before the period, volume and
atomic certificate; taking a maximum over the four entries preserves it.

The transported field and its local comparison are, by definition,
\[
 v_i(\nu)=\vartheta_i^{-2}P\mathcal C_{k_i,\nu}B_i(\nu),\qquad
 y_i(\nu)=P(\vartheta_i^{-2}X_i(\nu)).
\]
The column here is exactly the one defining that bias, so
\begin{equation}
 \|v_i(\nu)-y_i(\nu)\|_\infty\le A_{-1/4}(i).
 \label{eq:trace-weighted-normalized-bias}
\end{equation}
No bound for $\vartheta_i^{-2}$ is lost: it cancels the factor
$\vartheta_i^2$ in \eqref{eq:trace-weighted-paid-bias}.
Proposition~\ref{prop:ex-lattice-charge-bounds} and the payment accuracy
bound following it give
\[
 \|y_i(\nu)\|\le A_{1/6}(i),\qquad
 |\lambda(y_i(\nu))+\kappa|
 \le4\|\lambda\|A_{-1}(i)
       +(p_i^{-30}+p_i^{-100})/12.
\]
Since $A_{-1/4}\le A_{1/6}$, the triangle inequality now yields
\[
 \|\mathcal T v_i(\nu)\|\le2\|\mathcal T\|A_{1/6}(i),
\]
and
\[
 |\lambda(v_i(\nu))+\kappa|
 \le\|\lambda\|A_{-1/4}(i)+4\|\lambda\|A_{-1}(i)
             +(p_i^{-30}+p_i^{-100})/12.
\]
Thus \eqref{eq:ex-lattice-trace-finite-response-price} holds with
$C_T=2\|\mathcal T\|$, $C_1=\|\lambda\|$ and $C_2=4\|\lambda\|$.
The maximum norm on direction arrays is why no factor four occurs in
the transport error.

These fields really admit the fixed-level limit used above. Here is
the uniform bound needed before invoking compactness. Write $c_{j+1}$
for the stage column allowance and $L_j$ for the block side. The adjacent
calibrated-column inequality is
\[
 |\mathcal C_{j,\nu}b|_{\rm coeff}
 \le c_{j+1}L_j^d|\mathcal C_{j+1,\nu}b|_{\rm coeff}.
\]
At the reference level the column is $b$ itself. Induction over the
finitely many levels up to $k$ gives
\begin{equation}
 |b|_{\rm coeff}\le B_{d,k}|\mathcal C_{k,\nu}b|_{\rm coeff},
 \qquad B_{d,k}=\prod_{0\le j<k}c_{j+1}L_j^d.
 \label{eq:trace-weighted-column-inverse}
\end{equation}
The factors below the reference level only enlarge the bound, since they
are at least one. This constant depends on $d,k$ but not on the volume,
period or label. The finite-volume response record column has coefficient size at most two by
Lemma~\ref{lem:ex-trace-bill-column-size}, so its public label has norm at
most $2B_{d,k}$. A finite-dimensional closed ball is compact; the fixed
separation ultrafilter therefore gives a limit of these public labels.
Changing labels outside its large set has no effect. The joint
moving-label convergence defining the arch column sends this limit to the
limit of $\mathcal C_{k_i,\nu}B_i(\nu)$. Applying the fixed maps
$\vartheta_i^{-2}P$ and $\mathcal T$ gives precisely $q_i$, not a second
response chosen after taking the limit.

The same argument applies to all three reference profiles, not just the
one whose limit is the equation-of-motion vector. Write $P\in\{*,0,1\}$
for the stress profile and the two correction profiles. Denote their
finite-volume local probes by $c_i^P(\nu)$ and the corresponding
transported fields by $v_i^P(\nu)$. In the local radial chart, the
canonical probe has zero unit coordinate and its paid coefficient is
exactly $\vartheta_i^2c_i^P(\nu)$. The first assertion follows by evaluating
the probe at the zero field; the second is the definition of its weighted
local projection. Subsection~\ref{subsec:trace-reference-atomic-certificate} supplies
this full observable with charge $\vartheta_i^2A_{1/6}(i)$ and
attainment bound at most $p_i$. Applying
\eqref{eq:trace-weighted-paid-bias} to this atomic coefficient and then deleting the unit coordinate gives
\begin{equation}
 \|v_i^P(\nu)-c_i^P(\nu)\|\le A_{-1/4}(i).
 \label{eq:trace-weighted-reference-bias}
\end{equation}
Thus the reference uses its own probe and its own paid public label;
it is not compared to the chosen-response field $v_i(\nu)$.

Here is the precise finite-dimensional calibration needed for this
passage. Let $a^P\in\mathbb R^3$ be the coefficient triple of the
limiting profile and $V:\mathbb R^3\to E$ the fixed raw vector map.
The local probe classification proved in
Subsection~\ref{subsec:trace-reference-local-calibration} gives triples $a_i^P(\nu)$ with
\begin{equation}
 \mathcal T c_i^P(\nu)=V a_i^P(\nu),\qquad
 \rho(Va_i^P(\nu))=\rho(Va^P).
 \label{eq:trace-weighted-calibration-price}
\end{equation}
For the two elementary correction profiles the triples are exactly
$a^0,a^1$. For the stress profile the local stress probe tends to the
negative Maxwell tensor. The injective grade-four stress-coordinate map
of Lemma~\ref{lem:ex-trace-free-stress-basis} has a continuous left inverse,
so its coefficient triples tend to $a^*=(-1,1/4,-1)$. This convergence
is uniform in the local configurations on the late-level, sufficiently
large-period sets: otherwise choose a violating sequence with levels and
periods tending to infinity and apply the same local-probe limit and
left inverse to that sequence. Consequently, for every $\epsilon>0$,
\[
 \|a_i^P(\nu)-a^P\|\le\epsilon
\]
on an ultrafilter-large set for every sufficiently large $i$.

Put $U_i^P=\lim_{\nu\to\mathcal L}\mathcal T v_i^P(\nu)$ and
$v_P=Va^P$. By \eqref{eq:trace-weighted-reference-bias},
\[
 \|U_i^P-v_P\|
 \le\|\mathcal T\|A_{-1/4}(i)+\|V\|\epsilon
\]
for every sufficiently large $i$. First let $i\to\infty$, then
$\epsilon\downarrow0$, to obtain $U_i^P\to v_P$.
For the bound, use the exact identity
\eqref{eq:trace-weighted-calibration-price}, rather than the qualitative
convergence of the triples. It gives directly
\[
 |\rho(U_i^P)-\rho(v_P)|
 \le\|\rho\|\,\|\mathcal T\|A_{-1/4}(i),\qquad
 A_{1/5}(i)|\rho(U_i^P)-\rho(v_P)|
 \le\|\rho\|\,\|\mathcal T\|A_{-1/20}(i)\longrightarrow0.
\]
For $P=*$, the equation-of-motion identification is $v_*=-x$ and
$\rho(x)=0$; taking $w_i=U_i^*$ proves
\eqref{eq:ex-lattice-trace-reference-inputs-source}. The other two cases
are the weighted correction-profile estimates used by the stress transfer.
No rate of convergence of $a_i^*(\nu)$ to $a^*$ has been assumed.

The remaining chart-field input required by the general lemma is the exact
limit
\begin{equation}
 p_i\,s_i(m_i)\longrightarrow0.
 \label{eq:ex-lattice-trace-chart-field-decay}
\end{equation}
It is independent of the response-size and flux estimates above. It is the
conclusion of Lemma~\ref{lem:ex-trace-small-means} with exponent $1$,
applied to the present medium: in the notation of that lemma,
$z_i=p_i$, $\zeta_i=b_i$, and its medium $v_i$ is the present $m_i$.
The hypotheses of that lemma are the scale relations, the core
representation \eqref{eq:ex-trace-medium-core} of this medium, the
response and scaling identities, the local estimate for finitely many
detector fields, the detector lower bound
\eqref{eq:ex-trace-small-means-detector}, and the scalar rate
\eqref{eq:ex-trace-small-means-rate}. The detector lower bound is
essential. Bounds for the pairings of $s_i(m_i)$ with finitely many tested
fields give a bound for its norm only because those tests separate the
finite-dimensional response space with a lower bound uniform in $i$;
finite dimensionality alone does not provide this. All the hypotheses
listed are inputs to this step, and no claim that they follow from the
scalar rate calculation is made here.

For completeness the medium identity uses the same arch column as these
limits. Let $R$ denote signed averaging, $\mathcal E_k$ the normalized
core column of the arch, and $a_k,a_k^*$ the limiting public chosen and
stress labels. Signed equivariance of the common-trajectory columns gives
$R\mathcal E_k=\mathcal E_kR$, componentwise. Hence, with $b=b(k)>0$,
\[
 b^{-1}R\mathcal E_k a_k-R\mathcal E_k a_k^*
 =b^{-1}\mathcal E_k R(a_k-ba_k^*).
\]
The left side is exactly the medium already constructed, and the right
side is the core representation \eqref{eq:ex-trace-medium-core}, with
$x_i=R(a_{k_i}-b_i a_{k_i}^*)$. Both terms in this identity use the same
arch: replacing one by a separately chosen column limit would invalidate
it. The response identity in the small-means lemma is the defining
normalization of $s_i$. Central inversion acts as minus the identity on
these vector arrays, so its odd-grade scaling identity applies. Thus,
when the same core is inserted in the priced mission of
Proposition~\ref{prop:ex-trace-mission-price-transfer}, the resulting
field is the present $s_i(m_i)$. The detector and the bound used there
remain the substantive field estimates, rather than consequences of the
coefficient-size bound just proved.
